\begin{solapartat}
\figura[0.45]{sol_fig1.png}

\punts{0,25} Del gràfic $g = 0,6\ \mathrm{m/s^2}$.

La segona llei de Newton estableix que:

\punts{0,1} \[ \vec{F} = m_{S2}\vec{a} \]

I, la acceleració centrípeta de S2 és:

\punts{0,1} \[ a = \frac{v^2}{r} \]

La força gravitatòria s'expressa com:

\punts{0,4} \[ F = m_{S2}\,g \]

\punts{0,3} Per tant obtenim que: $\dfrac{v^2}{r} = g = 0,6\ \mathrm{m/s^2}$

\punts{0,1} I finalment: $v = \sqrt{rg} = 4,24\times 10^{6}\ \mathrm{m/s}$
\end{solapartat}

\begin{solapartat}
\punts{0,25} Segons la llei de gravitació universal, el mòdul de la força s'expressa com:
\[ F = G\,\frac{m_{S2} M_{A*}}{r^2} \]
\[ \text{i } F = m_{S2}\,g \]
Per tant obtenim que:

\punts{0,25} \[ g = G\,\frac{M_{A*}}{r^2} = 0,6\ \mathrm{m/s^2} \]

\punts{0,25} \[ M_{A*} = g\,\frac{r^2}{G} = 8,10\times 10^{36}\ \mathrm{kg}\,\frac{M_{Sol}}{1,99\times 10^{30}\ \mathrm{kg}} = 4,07\times 10^{6} M_{Sol} \]

\punts{0,25} El càlcul de la velocitat d'escapament es basa en el principi de conservació de l'energia mecànica. Per escapar cal que l'energia mecànica sigui igual a la de l'infinit que és 0: $E_m$(a una distància $r$ de Sagitari A*)${}= 0$
\[ E_m = E_c + E_p = 0 \Rightarrow \frac{1}{2} m v_{esc}^2 - G\,\frac{m M_{A*}}{r} = 0 \]

\punts{0,25}
\begin{align*}
\frac{1}{2} m v_{esc}^2 = G\,\frac{m M_{A*}}{r} \Rightarrow r &= \frac{2 G M_{A*}}{v_{esc}^2}\\
&= \frac{2\times 6,67\times 10^{-11}\cdot 8,10\times 10^{36}}{\left(3,00\times 10^{8}\right)^2} = 1,20\times 10^{10}\ \mathrm{m}
\end{align*}
\end{solapartat}
